% ======================================================================
\chapter[Волны де Бройля. Дифракция электронов. Соотношения неопределённостей]%
        {Волны де Бройля. Дифракция электронов. Соотношения неопределённостей}
\label{t:02}
% ======================================================================

\begin{qbox}
\textbf{Билет 2.} Волновые свойства массивных частиц — корпускулярно-волновой дуализм.
Опыты Девиссона--Джермера и Томсона по дифракции электронов. Длина волны де Бройля
нерелятивистской частицы. Соотношения неопределенностей (координата-импульс; энергия-время).
\end{qbox}

\section{Содержание билета}

\textbf{Подтемы:}
\begin{enumerate}
 \item гипотеза де Бройля и корпускулярно-волновой дуализм;
 \item длина волны де Бройля нерелятивистской частицы; электрон, ускоренный напряжением $U$;
 \item опыт Девиссона--Джермера (отражение от монокристалла) и опыт Томсона (прохождение через поликристаллическую фольгу);
 \item соотношение неопределённостей координата--импульс: формулировка, обоснование, примеры;
 \item соотношение неопределённостей энергия--время и его смысл.
\end{enumerate}

\textbf{Необходимые понятия:} фотон (\tref{t:01}); интерференция и дифракция волн;
условие Брэгга--Вульфа $2d\sin\vartheta=n\lambda$; фурье-разложение волнового пакета;
дисперсия (среднеквадратичное отклонение) случайной величины.

\section{Определения и физические принципы}

\begin{principle}{Гипотеза де Бройля (1924)}
С движением любой частицы с энергией $E$ и импульсом $\vect p$ связан волновой процесс
с частотой и волновым вектором
\begin{equation}\label{eq:2:dB}
 \omega=\frac{E}{\hbar},\qquad \vect k=\frac{\vect p}{\hbar},\qquad
 \lambda=\frac{2\pi}{k}=\frac{h}{p}.
\end{equation}
\end{principle}

\emph{Физический смысл.} Соотношения~\eqref{eq:2:dB} совпадают с соотношениями для фотона, но
применяются к частицам с ненулевой массой. Волна де Бройля — не классическая волна какого-либо
поля: её интенсивность определяет вероятность обнаружить частицу (статистическая трактовка
Борна, \tref{t:04}).

\begin{definition}[корпускулярно-волновой дуализм]
Свойство микрообъектов проявлять как корпускулярные (локализованная регистрация, сохранение
энергии и импульса в отдельном акте), так и волновые (интерференция, дифракция) свойства.
Какое из них проявляется, определяется условиями опыта; полное описание даёт квантовая механика.
\end{definition}

\begin{definition}[неопределённость]
Неопределённость величины $A$ в данном состоянии — среднеквадратичное отклонение результатов
её измерения от среднего по множеству одинаково приготовленных систем:
$\Delta A=\sqrt{\avg{(A-\avg A)^2}}$.
\end{definition}

\emph{Область применимости.} Волновые свойства существенны, когда длина волны де Бройля сравнима
с характерными размерами задачи (межатомные расстояния, размер атома, ширина щели). Для
макроскопических тел $\lambda$ ничтожно мала: пылинка массой $10^{-12}$~г со скоростью 1~см/с
имеет $\lambda=h/mv\approx7\cdot10^{-15}$~см, и волновые эффекты ненаблюдаемы.

\section{Основная теория}

\subsection{От фотона к частицам}

После опытов с фотоэффектом и эффектом Комптона стало ясно, что свет, считавшийся волной,
обладает корпускулярными свойствами. Л.~де Бройль предположил обратное: частицы вещества
обладают волновыми свойствами, причём связь между волновыми и корпускулярными характеристиками
та же, что у фотона. Аргумент: условие квантования Бора (\tref{t:03}) $L=n\hbar$ для круговой
орбиты эквивалентно требованию, чтобы на длине орбиты укладывалось целое число длин волн:
$2\pi r=n\lambda$, $\lambda=h/p$ $\Rightarrow$ $pr=n\hbar$.

\subsection{Длина волны де Бройля}

Для нерелятивистской частицы $p=mv=\sqrt{2mT}$, где $T$ — кинетическая энергия:
\begin{equation}\label{eq:2:lam}
  \lambda=\frac{h}{mv}=\frac{h}{\sqrt{2mT}}.
\end{equation}
Электрон, прошедший ускоряющую разность потенциалов $U$, имеет $T=eU$ и
\begin{equation}\label{eq:2:lamU}
  \lambda=\frac{h}{\sqrt{2\mel eU}}\approx\frac{12{,}26}{\sqrt{U[\text{В}]}}\ \text{Å}.
\end{equation}
При $U=100$~В $\lambda\approx1{,}23$~Å — порядок межатомных расстояний в кристаллах.
Поэтому кристалл служит дифракционной решёткой для электронов, как и для рентгеновских лучей.
Для частиц тепловых энергий $T\sim k_{\mathrm B}T_0$ ($T_0$ — температура) $\lambda\sim h/\sqrt{2mk_{\mathrm B}T_0}$:
для тепловых нейтронов при 300~К это $\approx 1{,}8$~Å, что используется в нейтронографии.

\begin{remark}[Релятивистская поправка]
При $eU\gtrsim0{,}01\,\mel c^2\approx5$~кэВ следует пользоваться точным выражением
$pc=\sqrt{T(T+2\mel c^2)}$, откуда
$\lambda=h\big/\sqrt{2\mel eU\,(1+eU/2\mel c^2)}$. Поправка $\approx eU/4\mel c^2$ к длине волны
важна в электронной микроскопии (100~кэВ — около 5\,\%).
\end{remark}

\subsection{Опыт Девиссона--Джермера (1927)}

Пучок электронов с энергией в десятки--сотни эВ падал нормально на грань монокристалла никеля;
детектор (цилиндр Фарадея) измерял ток отражённых электронов в зависимости от угла
(рис.~\ref{fig:2:dg}). При $U=54$~В под углом $\varphi=50^\circ$ к падающему пучку наблюдался
резкий максимум, который невозможно объяснить рассеянием частиц: он возникает только при
определённом напряжении.

\emph{Объяснение.} Поверхностный слой атомов — двумерная решётка с периодом $d=2{,}15$~Å
(расстояние между атомными рядами). Волны, рассеянные соседними рядами, усиливаются при
\begin{equation}\label{eq:2:grating}
  d\sin\varphi=n\lambda .
\end{equation}
При $n=1$: $\lambda=2{,}15\cdot\sin50^\circ\approx1{,}65$~Å. По формуле~\eqref{eq:2:lamU}
$\lambda=12{,}26/\sqrt{54}\approx1{,}67$~Å. Совпадение подтвердило формулу де Бройля
количественно. (Учёт преломления электронной волны на границе кристалла, где кинетическая энергия
электрона больше на величину внутреннего потенциала, устраняет оставшиеся небольшие расхождения
в положении максимумов высших порядков.)

\begin{figure}[htbp]
\centering
\begin{tikzpicture}[scale=0.95]
  % кристалл
  \foreach \x in {-2.4,-1.6,...,2.4} { \foreach \y in {0,-0.5,-1.0} { \fill[gray!60] (\x,\y) circle (2.2pt); } }
  \draw[gray] (-2.8,0.25) -- (2.8,0.25);
  \node[right, font=\small] at (2.6,-0.5) {кристалл Ni};
  % пушка
  \draw[thick] (-0.25,3.6) rectangle (0.25,4.4);
  \node[left, font=\small] at (-0.3,4.0) {электронная пушка};
  \draw[very thick, accent, ->] (0,3.55) -- (0,0.35);
  \node[right, font=\small, accent] at (0.05,2.2) {$e^-$, $eU$};
  % отражённый пучок
  \draw[very thick, keycol, ->] (0,0.35) -- ({0+3.2*sin(50)},{0.35+3.2*cos(50)});
  \draw[thick] ({3.35*sin(50)-0.25},{0.35+3.35*cos(50)}) rectangle ++(0.5,0.45);
  \node[right, font=\small] at ({3.35*sin(50)+0.3},{0.55+3.35*cos(50)}) {детектор};
  \draw[->] (0,1.5) arc[start angle=90, end angle=40, radius=1.15];
  \node at (0.6,1.75) {$\varphi$};
  % схема двух рядов
  \begin{scope}[xshift=7.2cm, yshift=0.4cm]
    \fill[gray!60] (0,0) circle (2.4pt); \fill[gray!60] (1.6,0) circle (2.4pt);
    \draw[<->] (0,-0.35) -- (1.6,-0.35) node[midway, below] {$d$};
    \draw[accent, ->] (0,2.2) -- (0,0.08); \draw[accent, ->] (1.6,2.2) -- (1.6,0.08);
    \draw[keycol, ->] (0,0) -- ({-2.2*sin(50)},{2.2*cos(50)});
    \draw[keycol, ->] (1.6,0) -- ({1.6-2.2*sin(50)},{2.2*cos(50)});
    \draw[dashed] (1.6,0) -- ({1.6*cos(50)*cos(50)},{1.6*cos(50)*sin(50)});
    \node[font=\small, align=left] at (0.2,-1.15) {разность хода $d\sin\varphi$};
  \end{scope}
\end{tikzpicture}
\caption{Опыт Девиссона--Джермера: схема установки (слева) и разность хода волн,
рассеянных соседними атомными рядами поверхности (справа).}
\label{fig:2:dg}
\end{figure}

\subsection{Опыт Томсона (1927)}

Дж.\,П.~Томсон пропускал пучок быстрых электронов (энергии $\sim$ 10--60~кэВ) через тонкую
($\sim10^{-6}$~см) поликристаллическую металлическую фольгу и регистрировал прошедшие электроны
на фотопластинке на расстоянии $L$ за фольгой. Наблюдалась система концентрических колец —
дебаеграмма, такая же, как при рентгеновской дифракции на поликристалле (метод Дебая--Шеррера).
Среди множества хаотически ориентированных микрокристаллов всегда есть такие, для которых
выполнено условие Брэгга--Вульфа
\begin{equation}\label{eq:2:bragg}
 2d\sin\vartheta=n\lambda ,
\end{equation}
где $d$ — межплоскостное расстояние, $\vartheta$ — угол скольжения. Отражённые лучи отклоняются
на угол $2\vartheta$ и образуют конус, т.\,е. кольцо радиуса $r\approx L\tg2\vartheta\approx2L\vartheta$.
При малых углах $\lambda\approx 2d\vartheta/n=d\,r/(nL)$, т.\,е. $r\propto\lambda\propto1/\sqrt U$, что
и наблюдалось. То, что кольца образованы именно электронами, а не вторичным рентгеновским
излучением, проверялось отклонением картины магнитным полем: кольца смещались целиком.

\paragraph{Дифракция одиночных частиц.} В опытах Л.~Бибермана, Н.~Сушкина и В.~Фабриканта (1949)
интенсивность пучка была столь мала, что электроны проходили прибор поодиночке. При длительной
экспозиции получалась та же дифракционная картина. Следовательно, волновые свойства присущи
каждому отдельному электрону, а не возникают от взаимодействия многих частиц. Аналогичные
результаты получены для нейтронов, атомов и молекул (вплоть до фуллеренов C$_{60}$).

\subsection{Соотношение неопределённостей координата--импульс}

\begin{principle}{Соотношение неопределённостей Гейзенберга}
В любом состоянии частицы неопределённости координаты и соответствующей проекции импульса
удовлетворяют неравенству
\begin{equation}\label{eq:2:HUP}
  \Delta x\,\Delta p_x\ge\frac{\hbar}{2},\qquad
  \Delta y\,\Delta p_y\ge\frac{\hbar}{2},\qquad
  \Delta z\,\Delta p_z\ge\frac{\hbar}{2}.
\end{equation}
Для «разноимённых» пар ($x$ и $p_y$) ограничений нет.
\end{principle}

\emph{Происхождение.} Частица с точно определённым импульсом описывается монохроматической
волной $\ee^{\ii p x/\hbar}$, распределённой по всему пространству, — её координата полностью
неопределённа. Чтобы локализовать частицу в области $\Delta x$, нужно сложить волны с разными
$k$ (волновой пакет). Из теории фурье-разложения известно, что ширины пакета в координатном и
$k$-пространствах связаны: $\Delta x\,\Delta k\gtrsim1$ (точнее, $\ge1/2$ для среднеквадратичных
ширин). С учётом $p=\hbar k$ это и есть~\eqref{eq:2:HUP}. Строгое доказательство для любых
операторов дано в \tref{t:05}; знак равенства достигается на гауссовом пакете
(вывод~\ref{der:2:gauss}).

\emph{Что соотношение НЕ означает.} Оно не связано с несовершенством приборов: это свойство
самого состояния. Частица не имеет одновременно определённых $x$ и $p_x$, поэтому понятие
траектории для микрочастиц теряет смысл. Классическое описание применимо, когда $\Delta x$ и
$\Delta p$, допускаемые~\eqref{eq:2:HUP}, пренебрежимо малы по сравнению с масштабами задачи.

\paragraph{Мысленный эксперимент: дифракция на щели.} Пусть частица с импульсом $p$ проходит
щель ширины $b$. Сразу после щели координата известна с точностью $\Delta x\sim b$. Из-за дифракции
частица отклоняется в пределах центрального максимума: $\sin\theta_1=\lambda/b$, т.\,е. приобретает
поперечный импульс $\Delta p_x\sim p\sin\theta_1=p\lambda/b=h/b$. Отсюда $\Delta x\,\Delta p_x\sim h$:
чем уже щель, тем шире дифракционная картина.

\subsection{Соотношение неопределённостей энергия--время}

\begin{principle}{Соотношение энергия--время}
\begin{equation}\label{eq:2:Et}
 \Delta E\,\Delta t\gtrsim\hbar .
\end{equation}
\end{principle}

Это соотношение имеет иной смысл, чем~\eqref{eq:2:HUP}: время в квантовой механике — параметр,
а не наблюдаемая с оператором. Существуют следующие корректные трактовки.
\begin{enumerate}
 \item \textbf{Время жизни и ширина уровня.} Состояние, существующее время $\tau$ (например,
       возбуждённое состояние атома), не имеет точно определённой энергии: его энергия размыта
       на ширину $\Gamma\approx\hbar/\tau$. Это проявляется как естественная ширина спектральной
       линии.
 \item \textbf{Длительность измерения.} Чтобы измерить энергию с точностью $\Delta E$, требуется
       время $\Delta t\gtrsim\hbar/\Delta E$ (волновой цуг длительности $\Delta t$ имеет ширину спектра
       $\Delta\omega\sim1/\Delta t$).
 \item \textbf{Характерное время изменения состояния.} Если $\Delta t$ — время, за которое средние
       значения величин заметно меняются, то разброс энергии в состоянии $\Delta E\gtrsim\hbar/\Delta t$
       (соотношение Мандельштама--Тамма, вывод опущен). Стационарное состояние ($\Delta E=0$) не
       меняется со временем вовсе.
\end{enumerate}

\section{Математические выводы}

\begin{derivation}[label=der:2:lamU]{длина волны электрона, ускоренного напряжением $U$}
Электрон, ускоренный из состояния покоя, приобретает кинетическую энергию $T=eU$. При
$eU\ll\mel c^2$ импульс $p=\sqrt{2\mel T}=\sqrt{2\mel eU}$, и
\[
 \lambda=\frac{h}{\sqrt{2\mel eU}}.
\]
Численно, используя $hc=12\,398$~эВ$\cdot$Å и $\mel c^2=0{,}511\cdot10^6$~эВ:
\[
 \lambda=\frac{hc}{\sqrt{2\,\mel c^2\,eU}}
 =\frac{12\,398\ \text{эВ}\cdot\text{Å}}{\sqrt{2\cdot0{,}511\cdot10^6\ \text{эВ}\cdot U[\text{В}]\ \text{эВ}}}
 =\frac{12\,398}{1011\sqrt{U[\text{В}]}}\ \text{Å}\approx\frac{12{,}26}{\sqrt{U[\text{В}]}}\ \text{Å}.
\]
Размерность: $[h/\sqrt{mT}]=\dfrac{\text{эрг}\cdot\text{с}}{\sqrt{\text{г}\cdot\text{эрг}}}=\dfrac{\text{г}\,\text{см}^2/\text{с}}{\text{г}\,\text{см}/\text{с}}=\text{см}$.
\end{derivation}

\begin{derivation}[label=der:2:packet]{оценка $\Delta x\,\Delta k\gtrsim 1$ для волнового пакета}
Рассмотрим пакет, составленный из волн с волновыми числами в интервале
$k_0-\Delta k/2<k<k_0+\Delta k/2$ с одинаковой амплитудой (момент $t=0$):
\[
 \psi(x)=\int_{k_0-\Delta k/2}^{k_0+\Delta k/2}\ee^{\ii kx}\,\dd k
 =\ee^{\ii k_0x}\,\frac{\ee^{\ii\Delta k\,x/2}-\ee^{-\ii\Delta k\,x/2}}{\ii x}
 =\ee^{\ii k_0x}\,\Delta k\,\frac{\sin(\Delta k\,x/2)}{\Delta k\,x/2}.
\]
Модуль огибающей максимален при $x=0$ и первый раз обращается в нуль при $\Delta k\,x/2=\pi$,
т.\,е. при $x=2\pi/\Delta k$. Принимая за ширину пакета половину расстояния между нулями,
$\Delta x\approx2\pi/\Delta k$, получаем $\Delta x\,\Delta k\approx2\pi$, а с $p=\hbar k$ —
$\Delta x\,\Delta p_x\approx h$. Конкретное число зависит от определения ширины, но
произведение не может быть сделано сколь угодно малым: сужение пакета в одном пространстве
неизбежно уширяет его в другом.
\end{derivation}

\begin{derivation}[label=der:2:gauss]{гауссов пакет: $\Delta x\,\Delta p_x=\hbar/2$}
Пусть $\psi(x)=(2\pi\sigma^2)^{-1/4}\exp\!\left(-\dfrac{x^2}{4\sigma^2}+\ii k_0x\right)$. Тогда
$|\psi|^2=(2\pi\sigma^2)^{-1/2}\ee^{-x^2/2\sigma^2}$ — нормальное распределение с дисперсией
$\sigma^2$, т.\,е. $\avg{x}=0$, $\Delta x=\sigma$.

Амплитуда распределения по $k$ — фурье-образ:
\[
 a(k)=\frac{1}{\sqrt{2\pi}}\int_{-\infty}^{\infty}\psi(x)\ee^{-\ii kx}\dd x
 =\frac{(2\pi\sigma^2)^{-1/4}}{\sqrt{2\pi}}\int_{-\infty}^{\infty}
   \exp\!\Bigl[-\frac{x^2}{4\sigma^2}-\ii(k-k_0)x\Bigr]\dd x .
\]
Используем интеграл Пуассона $\int_{-\infty}^{\infty}\ee^{-\alpha x^2-\ii\beta x}\dd x=\sqrt{\pi/\alpha}\,\ee^{-\beta^2/4\alpha}$
при $\alpha=1/4\sigma^2$, $\beta=k-k_0$:
\[
 a(k)\propto\exp\bigl[-\sigma^2(k-k_0)^2\bigr],\qquad
 |a(k)|^2\propto\exp\Bigl[-\frac{(k-k_0)^2}{2\,(1/2\sigma)^2}\Bigr].
\]
Это нормальное распределение с дисперсией $(1/2\sigma)^2$: $\Delta k=1/(2\sigma)$. Следовательно,
\[
 \Delta x\,\Delta k=\sigma\cdot\frac{1}{2\sigma}=\frac12,\qquad
 \Delta x\,\Delta p_x=\frac{\hbar}{2}.
\]
Гауссов пакет — состояние с минимально возможным произведением неопределённостей.
\end{derivation}

\begin{derivation}[label=der:2:atom]{оценка размера и энергии атома водорода из соотношения неопределённостей}
Электрон локализован в области размером $r$, поэтому $p\sim\Delta p\sim\hbar/r$. Энергия
\[
 E(r)\approx\frac{p^2}{2\mel}-\frac{e^2}{r}\approx\frac{\hbar^2}{2\mel r^2}-\frac{e^2}{r}.
\]
Минимизируем: $\dfrac{\dd E}{\dd r}=-\dfrac{\hbar^2}{\mel r^3}+\dfrac{e^2}{r^2}=0$
$\Rightarrow$ $r_0=\dfrac{\hbar^2}{\mel e^2}=\aB$, и
$E(r_0)=\dfrac{\hbar^2}{2\mel}\dfrac{\mel^2e^4}{\hbar^4}-\dfrac{\mel e^4}{\hbar^2}=-\dfrac{\mel e^4}{2\hbar^2}=-13{,}6\eV$.
Совпадение с точным результатом (\tref{t:10}) отчасти случайно (оценка верна по порядку
величины), но рассуждение показывает главное: \emph{атом устойчив потому, что сжатие электронного
облака увеличивает кинетическую энергию быстрее, чем уменьшает потенциальную.}
\end{derivation}

\begin{derivation}[label=der:2:nucleus]{почему электрон не может находиться в ядре}
Если бы электрон был локализован в области ядра $\Delta x\sim R\sim5\cdot10^{-13}$~см, то
$p\gtrsim\hbar/R$, и $pc\gtrsim\hbar c/R=197\MeV\cdot\text{фм}/5\fm\approx40$~МэВ $\gg\mel c^2$.
Электрон был бы ультрарелятивистским с энергией $\sim40$~МэВ, тогда как энергии электронов
бета-распада — порядка 1~МэВ, а глубина ядерного потенциала не удерживает электрон (на него
не действуют ядерные силы). Значит, электроны бета-распада рождаются в момент распада
(\tref{t:20}).
\end{derivation}

\section{Итоговые формулы}

\keyformula{\lambda=\frac{h}{p}=\frac{h}{mv},\qquad \vect k=\frac{\vect p}{\hbar},\qquad \omega=\frac{E}{\hbar}}
 {$p$ — импульс частицы, $m$ — масса, $v$ — скорость.}
 {любая частица; $p=mv$ — нерелятивистский случай.}
 {длина волны де Бройля; оценка роли волновых эффектов.}

\keyformula{\lambda=\frac{h}{\sqrt{2\mel eU}}\approx\frac{12{,}26}{\sqrt{U[\text{В}]}}\ \text{Å}}
 {$U$ — ускоряющее напряжение.}
 {электрон, $eU\ll\mel c^2$ (практически до нескольких кэВ).}
 {длина волны в опытах по дифракции электронов.}

\keyformula{d\sin\varphi=n\lambda\ \ \text{(Девиссон--Джермер)},\qquad 2d\sin\vartheta=n\lambda\ \ \text{(Брэгг--Вульф)}}
 {$d$ — период поверхностной решётки / межплоскостное расстояние, $\varphi$ — угол отражения от нормали, $\vartheta$ — угол скольжения, $n$ — порядок.}
 {однократное (кинематическое) рассеяние; без учёта преломления на поверхности.}
 {связь положения дифракционных максимумов с $\lambda$.}

\keyformula{\Delta x\,\Delta p_x\ge\frac{\hbar}{2},\qquad \Delta E\,\Delta t\gtrsim\hbar,\qquad \Gamma\tau\approx\hbar}
 {$\Delta$ — среднеквадратичные отклонения; $\Gamma$ — ширина уровня, $\tau$ — время жизни.}
 {первое — строгое неравенство для любого состояния; второе — оценочное, $t$ — не наблюдаемая.}
 {оценки размеров и энергий связанных состояний, ширины линий, времён жизни.}

\section{Физическая интерпретация}

\textbf{Масштабы.} $\lambda\propto1/\sqrt{mT}$: при одинаковой энергии тяжёлые частицы имеют
меньшую длину волны ($\lambda_p/\lambda_e=\sqrt{\mel/m_p}\approx1/43$). Электроны с энергией
100~эВ и рентгеновские фотоны с энергией $\sim10$~кэВ имеют длины волн порядка ангстрема —
отсюда электронография и рентгеноструктурный анализ. Высокая разрешающая способность электронного
микроскопа обусловлена малостью $\lambda$ при энергиях 100--300~кэВ.

\textbf{Классический предел.} Классическая механика применима, если $\lambda$ мала по сравнению
с масштабом изменения потенциала, а неопределённости, следующие из~\eqref{eq:2:HUP},
пренебрежимы: например, для электрона в кинескопе ($\Delta x\sim10^{-3}$~см) $\Delta v\sim\hbar/(\mel\Delta x)\sim10^3$~см/с
при $v\sim10^9$~см/с — траекторию можно использовать. Для электрона в атоме $\Delta v\sim v$ —
понятие орбиты теряет смысл.

\textbf{Энергия нулевых колебаний.} Частица, запертая в области размера $a$, не может иметь
нулевую кинетическую энергию: $E_{\min}\sim\hbar^2/(ma^2)$. Это объясняет конечную энергию
основного состояния в яме (\tref{t:06}) и осцилляторе (\tref{t:08}).

\textbf{Ширина линий.} Время жизни возбуждённого атома $\tau\sim10^{-8}$~с даёт $\Gamma\sim\hbar/\tau\sim
7\cdot10^{-8}$~эВ — естественная ширина линии. Короткоживущие адроны ($\tau\sim10^{-23}$~с) имеют
ширины $\sim100$~МэВ, и их «масса» определяется лишь с такой точностью (\tref{t:15}).

\section{Что необходимо знать на экзамене}

\begin{examcheck}
\checkhead{Определения}
\begin{itemize}
 \item гипотеза де Бройля: $\lambda=h/p$, $\omega=E/\hbar$; смысл корпускулярно-волнового дуализма;
 \item неопределённость как среднеквадратичное отклонение.
\end{itemize}
\checkhead{Формулы}
\begin{itemize}
 \item $\lambda=h/\sqrt{2mT}$, $\lambda\approx12{,}26/\sqrt{U}$~Å;
 \item $d\sin\varphi=n\lambda$, $2d\sin\vartheta=n\lambda$;
 \item $\Delta x\,\Delta p_x\ge\hbar/2$, $\Delta E\,\Delta t\gtrsim\hbar$.
\end{itemize}
\checkhead{Выводы и оценки}
\begin{itemize}
 \item численная обработка опыта Девиссона--Джермера: 54~В, $50^\circ$, $d=2{,}15$~Å $\Rightarrow$ $\lambda\approx1{,}65$~Å;
 \item оценка $\Delta x\,\Delta k\sim2\pi$ для пакета; дифракция на щели;
 \item оценка размера и энергии атома водорода; невозможность электрона в ядре.
\end{itemize}
\checkhead{Интерпретации}
\begin{itemize}
 \item чем различаются опыты Девиссона--Джермера и Томсона; как доказано, что кольца дают электроны;
 \item дифракция одиночных электронов — волновые свойства присущи каждой частице;
 \item смысл соотношения энергия--время (время жизни и ширина уровня).
\end{itemize}
\end{examcheck}
